\documentclass[t]{beamer}
\usepackage{pdfpages}
\usepackage{animate}
\usepackage{mathtools}
\usepackage{tikz}

\usetheme{Montpellier}
% or ...
% \usecolortheme{dove}
\useoutertheme{infolines}

% \setbeamercovered{transparent}
% or whatever (possibly just delete it)


\usepackage[english]{babel}
% or whatever

\usepackage[utf8]{inputenc}
% or whatever

\usepackage[T1]{fontenc}
% Or whatever. Note that the encoding and the font should match. If T1
% does not look nice, try deleting the line with the fontenc.


\title[3D straight-line drawings of planar graphs]{Three-dimensional straight-line drawings of planar graphs} % (optional, use only with long paper titles)

% \subtitle
% {Presentation Subtitle} % (optional)

\author
{Tobias Feigenwinter}

\institute
{ETH Zürich}
% - Use the \inst command only if there are several affiliations.
% - Keep it simple, no one is interested in your street address.

\date % (optional)
{2023/04/21}

\subject{Talks}
% This is only inserted into the PDF information catalog. Can be left
% out. 



% If you have a file called "university-logo-filename.xxx", where xxx
% is a graphic format that can be processed by latex or pdflatex,
% resp., then you can add a logo as follows:

% \pgfdeclareimage[height=0.5cm]{university-logo}{university-logo-filename}
% \logo{\pgfuseimage{university-logo}}



% Delete this, if you do not want the table of contents to pop up at
% the beginning of each subsection:

% If you wish to uncover everything in a step-wise fashion, uncomment
% the following command: 

%\beamerdefaultoverlayspecification{<+->}


\newcommand\still[2][1]{
    \includegraphics[page=#1]{#2.pdf}
}

\newcommand\footcite[1]{%
  \begingroup
  \renewcommand\thefootnote{}\footnote{#1}%
  \addtocounter{footnote}{-1}%
  \endgroup
}

\newtheorem{task}{Task}
\newtheorem{question}{Question}
\newtheorem{claim}{Claim}
\begin{document}
\begin{frame}
    \titlepage
\end{frame}

\section{Motivation for 3D drawings}
\begin{frame}{Why 3D drawings?}{}
    \only<+>{\relax}\begin{itemize}
	\item<+->
	    ``Nice'' drawings

	    \only<+> {
		\still{static/icosahedralgraph}
	    }
	    \only<+> {
		\still{static/cube}
	    }
	    \only<+> {
		\still{anim/icosahedron}
	    }
	    \only<+> {
		\animategraphics[autoplay, loop]{20}{anim/icosahedron}{}{}
	    }
	\item<+-> Visualization of data with spacial information
	    \only<+> {
		\still{anim/rand3d}
	    }
	    \only<+> {
		\animategraphics[autoplay, loop]{20}{anim/rand3d}{}{}
	    }
	\item<+->
	    User Interfaces 
	    \only<+>{
		\includegraphics[height=0.55\paperheight, page=5, trim={0 17cm 0 1.5cm}, clip]{figures/Conetrees.5.pdf}
		\footcite{Figure from: George, Jock and Stuart (1991): \textit{Cone Trees: Animated 3D Visualizations of Hierarchical Information}}
	    }
	    \only<+->{
		\begin{itemize}
		    \item Virtual Reality?
		\end{itemize}
	}
    \end{itemize}
\end{frame}

\begin{frame}{Main result}
    \footcite{Theorem from: Dujmović, Joret, Micek, Morin, Ueckerdt and Wood (2020): \textit{Planar Graphs Have Bounded Queue Number}}
    \begin{theorem}[Dujmović et al. 2020]
	\uncover<+->{Every planar graph with $n$ vertices has a 3-dimensional crossing-free straight-line drawing on the integer grid}
	\uncover<+->{with $\mathcal O(n)$ Volume.}
    \end{theorem}
    \begin{center}
	\animategraphics[autoplay, loop]{20}{anim/gridgraphexample}{}{}
    \end{center}
\end{frame}

\begin{frame}{Why is this interesting?}
\end{frame}

\begin{frame}{The Moment Curve}
    \vspace{-0.5cm}
    \begin{columns}
	\begin{column}{0.5\textwidth}
		\begin{itemize}
		    \item<+-> Place $i$-th vertex at $(i, i^2, i^3)$
		    \item<+-> Four vertices are coplanar iff this matrix has determinant zero:
			\begin{equation*}
			    M = \begin{pmatrix}
				\uncover<+->{1}&i_0&i_0^2&i_0^3\\
				\uncover<.->{1}&i_1&i_1^2&i_1^3\\
				\uncover<.->{1}&i_2&i_2^2&i_2^3\\
				\uncover<.->{1}&i_3&i_3^2&i_3^3\\
			    \end{pmatrix}
			\end{equation*}
		    \item<+-> This is a Vandermonde Matrix \rlap{with}
			\[\det(M) = \prod_{0 \le k < l < 4} \alt<+->{\underbrace{(i_l - i_k)}_{\neq 0}}{(i_l-i_k)} \uncover<+->{\neq 0}\]
		    \item<+-> No four vertices coplanar \linebreak $\implies$ no two edges cross
		\end{itemize}
	\end{column}
	    \begin{column}{0.5\textwidth}
		\begin{center}
		    \animategraphics[autoplay, loop]{20}{anim/momentcurve}{}{}
		\end{center}
	    \end{column}
    \end{columns}
\end{frame}

\begin{frame}{The Modular Moment Curve}
    \vspace{-0.5cm}
    \begin{columns}
	\begin{column}{0.5\textwidth}
		\begin{itemize}
		    \item <+-> Let $p$ prime, $n\le p<2n$
		    \item
			Place $i$-th vertex at \rlap{$(i, i^2 \uncover<+->{\alert{\mod p}}, i^3 \uncover<.->{\alert{\mod p}})$} \\
		    \item Four vertices are coplanar iff this matrix has determinant zero:
			\begin{equation*}
			    M = \begin{pmatrix}
				{1}&i_0&i_0^2\uncover<+->{\alert{\mod p}}&i_0^3\uncover<.->{\alert{\mod p}}\\
				{1}&i_1&i_1^2\uncover<.->{\alert{\mod p}}&i_1^3\uncover<.->{\alert{\mod p}}\\
				{1}&i_2&i_2^2\uncover<.->{\alert{\mod p}}&i_2^3\uncover<.->{\alert{\mod p}}\\
				{1}&i_3&i_3^2\uncover<.->{\alert{\mod p}}&i_3^3\uncover<.->{\alert{\mod p}}\\
			    \end{pmatrix}
			\end{equation*}
		    \item 
			\(\det(M) \alt<+->{\alert{\equiv}}{=} \prod_{0 \le k < l < 4} \underbrace{(i_l - i_k)}_{\alt<+->{\alert{\not\equiv}}{\neq}0} \alt<+->{\alert{\not\equiv}}{\neq}\mathrlap{0 \uncover<.(-2)->{\alert{\pmod p}}}\)
		    \item No four vertices coplanar \linebreak $\implies$ no two edges cross
		\end{itemize}
	\end{column}
	    \begin{column}{0.5\textwidth}
		\begin{center}
		    \only<-6>{
		    \animategraphics[autoplay, loop]{20}{anim/modularmomentcurve}{}{}
		}\only<7->{
		    \animategraphics[autoplay, loop]{20}{anim/modularmomentcompletegraph}{}{}
		}
		\end{center}
	    \end{column}
    \end{columns}
\end{frame}

\begin{frame}{Upper bound for general graphs}
    \begin{theorem}[Cohen et al. 1996]
	Every graph on $n$ vertices has a $\mathcal O(n^3)$ 3D crossing-free straight-line drawing.
    \end{theorem}
\end{frame}

\begin{frame}{Can we do better?}
    \begin{columns}
	\begin{column}{0.4\textwidth}
	    \pause We try drawing a complete graph
	    \begin{itemize}
		    \pause
		\item Cut grid into planes \pause
		\item At most four vertices per plane: $K_5$ doesn't fit\pause
		\item At least $\lceil\frac{n}{4}\rceil$ planes in each dimension\pause
		\item At least $\lceil\frac{n}{4}\rceil^{3} \in \Omega(n^3)$ Volume
	    \end{itemize}
	\end{column}
	\begin{column}{0.6\textwidth}
	    \only<*>{
		\begin{center}
		    \animategraphics[autoplay, loop]{20}{anim/upperboundgrid}{}{}
		\end{center}
	    }
	\end{column}
    \end{columns}
\end{frame}

\section{Proof of the main result}

\begin{frame}{Outline of the proof}
    \foreach \i in {1, ..., 3} {
	\only<+>{
	    \begin{center}
		\still[\i]{static/outline}
	    \end{center}
	}
    }
\end{frame}

\subsection{Track Layouts to 3D drawings}

\begin{frame}{Track Layouts}
    \foreach \i in {1, ..., 4} {
	\only<+>{
	    \begin{center}
		\still[\i]{static/tracklayout}
	    \end{center}
	}
    }
\end{frame}

\begin{frame}{The Modular Moment Curve with tracks}
    \vspace{-0.5cm}
    \begin{columns}
	\begin{column}{0.5\textwidth}
		\begin{itemize}
		    \item <+-> Let $p$ prime, $t\le p<2t$
		    \item
			Place $j$-th vertex of $i$-th track at \rlap{$(i, i^2 \mod p, i^3 \mod p \alert{+jp})$} \\
		\end{itemize}
	\end{column}
	    \begin{column}{0.5\textwidth}
		\begin{center}
		    \animategraphics[autoplay, loop]{20}{anim/trackmomentcurve}{}{}
		\end{center}
	    \end{column}
    \end{columns}
\end{frame}

\begin{frame}{The Modular Moment Curve with tracks}
    \vspace{-0.5cm}
    \begin{columns}
	\begin{column}{0.5\textwidth}
		\begin{itemize}
		    \item <+-> Let $p$ prime, $t\le p<2t$
		    \item
			Place $j$-th vertex of $i$-th track at \rlap{$(i, i^2 \mod p, i^3 \mod p \alert{+jp})$} \\
			\end{itemize}
			Case 1: All i's different
			\begin{itemize}
		    \item Four vertices are coplanar iff this matrix has determinant zero:
			\begin{equation*}
			    M = \begin{pmatrix}
				{1}&i_0&i_0^2\mod p&i_0^3\mod p \alert{+j_0p}\\
				{1}&i_1&i_1^2\mod p&i_1^3\mod p \alert{+j_1p}\\
				{1}&i_2&i_2^2\mod p&i_2^3\mod p \alert{+j_2p}\\
				{1}&i_3&i_3^2\mod p&i_3^3\mod p \alert{+j_3p}\\
			    \end{pmatrix}
			\end{equation*}
		    \item 
			\(\det(M) \equiv \mathrlap{\prod_{0 \le k < l < 4} \underbrace{(i_l - i_k)}_{\not\equiv 0} \not\equiv0 \pmod p}\)
		    \item No four vertices coplanar \linebreak $\implies$ no two edges cross
		\end{itemize}
	\end{column}
	    \begin{column}{0.5\textwidth}
		\begin{center}
		    \animategraphics[autoplay, loop]{20}{anim/trackmomentcurve}{}{}
		\end{center}
	    \end{column}
    \end{columns}
\end{frame}

\begin{frame}{The Modular Moment Curve with tracks}
    \vspace{-0.5cm}
    \begin{columns}
	\begin{column}{0.5\textwidth}
		\begin{itemize}
		    \item <+-> Let $p$ prime, $t\le p<2t$
		    \item
			Place $j$-th vertex of $i$-th track at \rlap{$(i, i^2 \mod p, i^3 \mod p \alert{+jp})$} \\
			\end{itemize}
			Case 2: \rlap{Two i's the same, WLOG $i_0=i_1$}
			\begin{itemize}
		    \item Four vertices are coplanar iff this matrix has determinant zero:
			\begin{equation*}
			    M = \begin{pmatrix}
				{1}&i_1&i_1^2\mod p&i_1^3\mod p \alert{+j_0p}\\
				{1}&i_1&i_1^2\mod p&i_1^3\mod p \alert{+j_1p}\\
				{1}&i_2&i_2^2\mod p&i_2^3\mod p \alert{+j_2p}\\
				{1}&i_3&i_3^2\mod p&i_3^3\mod p \alert{+j_3p}\\
			    \end{pmatrix}
			\end{equation*}
		    \item <+-> 
			\(\det(M) = \underbrace{(j_1p-j_0p)}_{\neq 0}\cdot\det
			    \underbrace{
				\begin{pmatrix}
				    {1}&i_1&i_1^2\mod p\\
				    {1}&i_2&i_2^2\mod p\\
				    {1}&i_3&i_3^2\mod p\\
				\end{pmatrix}
			    }_{\not\equiv0} \mathrlap{\neq 0}
			\)
		\end{itemize}
	\end{column}
	    \begin{column}{0.5\textwidth}
		\begin{center}
		    \animategraphics[autoplay, loop]{20}{anim/trackmomentcurve}{}{}
		\end{center}
	    \end{column}
    \end{columns}
\end{frame}

\begin{frame}{The Modular Moment Curve with tracks}
    \vspace{-0.5cm}
    \begin{columns}
	\begin{column}{0.5\textwidth}
		\begin{itemize}
		    \item <+-> Let $p$ prime, $t\le p<2t$
		    \item
			Place $j$-th vertex of $i$-th track at \rlap{$(i, i^2 \mod p, i^3 \mod p \alert{+jp})$} \\
			\end{itemize}
			Case 3: Two pairs of $i$s are the same, e.g. $i_0 = i_2$ and $i_1 = i_3$
			\begin{itemize}
			    \item<+-> Edges are between the same track
			    \item<+-> Crossing in the 3D drawing would also be crossing in track layout
		\end{itemize}
	\end{column}
	    \begin{column}{0.5\textwidth}
		\begin{center}
		    \animategraphics[autoplay, loop]{20}{anim/trackmomentcurve-xcr}{}{}
		\end{center}
	    \end{column}
    \end{columns}
\end{frame}

\begin{frame}{The Modular Moment Curve with tracks}
    \vspace{-0.5cm}
    \begin{columns}
	\begin{column}{0.5\textwidth}
		\begin{itemize}
		    \item <+-> Let $p$ prime, $t\le p<2t$
		    \item
			Place $j$-th vertex of $i$-th track at \rlap{$(i, i^2 \mod p, i^3 \mod p \alert{+jp})$} \\
			\end{itemize}
			Case 4: Three or more $i$s are the same
			\begin{itemize}
			    \item<+-> Impossible: \rlap{requires intertrack edges}
			\end{itemize}
	\end{column}
	    \begin{column}{0.5\textwidth}
		\begin{center}
		    \animategraphics[autoplay, loop]{20}{anim/trackmomentcurve-ite}{}{}
		\end{center}
	    \end{column}
    \end{columns}
\end{frame}

\begin{frame}{The Modular Moment Curve with tracks}
    \vspace{-0.5cm}
    \begin{columns}
	\begin{column}{0.5\textwidth}
	    \begin{theorem}[Dujmović et al. 2005]
		Every graph on $n$ vertices with track number $t$ has a 3-dimensional crossing-free straight-line drawing on the integer grid with $\mathcal O(t^3n)$ volume. 
	    \end{theorem}
	\end{column}
	    \begin{column}{0.5\textwidth}
		\begin{center}
		    \animategraphics[autoplay, loop]{20}{anim/trackmomentcurve}{}{}
		\end{center}
	    \end{column}
    \end{columns}
\end{frame}

\begin{frame}{Outline of the proof}
    \foreach \i in {4, ..., 5} {
	\only<+>{
	    \begin{center}
		\still[\i]{static/outline}
	    \end{center}
	}
    }
\end{frame}

\begin{frame}
    \begin{definition}[Acyclic coloring]
	An acyclic coloring is a proper coloring such that every bichromatic subgraph is a forest\uncover<6->{, or equivalently, such that every cycle receives at least three colors. }
    \end{definition}
    \uncover<7->{
    \begin{definition}[Acyclic Chromatic Number]
	The acyclic chromatic number $\chi_a(G)$ is the smallest number of colors across all acyclic colorings of $G$.
    \end{definition}
}
    \foreach\i in {1, ..., 5}{
	\only<+>{
	    \begin{center}
		\still[\i]{static/acycliccoloring}
	    \end{center}
	}
    }
	\only<+->{
	    \begin{center}
		\still[5]{static/acycliccoloring}
	    \end{center}
	}
\end{frame}

\begin{frame}{The Refinement Lemma}
	    \vspace{-0.5cm}
    \begin{columns}
	\begin{column}{0.5\textwidth}
	    \begin{lemma}[Alon and Marshall 1998]
		\only<+->{
		Given a graph with}
		\begin{itemize}
		    \item<+-> an acyclic {\only<3->{\relax}$c$}-coloring
		    \item<+-> a (not necessarily proper) {\only<3->{\relax}$k$}-edge-coloring,
		\end{itemize}
		\only<+->{
		    there is a proper $ck^{c-1}$-vertex-coloring of $G$ such that
		}
		\begin{itemize}
		    \item<+-> every new vertex color class is contained in one old vertex color class, \rlap{and}
		    \item<+-> the edges between any two vertex color classes are monochromatic. 
		\end{itemize}
	    \end{lemma}
	\end{column}
	\begin{column}{0.5\textwidth}%
		\only<-6>{%
		    \begin{center}%
			\still[1]{static/refinementlemma}%
		    \end{center}%
		}%
		\only<7->{%
		    \begin{center}%
			\still[2]{static/refinementlemma}%
		    \end{center}%
		}%
	\end{column}%
    \end{columns}
\end{frame}

\begin{frame}
    \begin{theorem}[Dujmović et al. 2005]
	Every graph with queue number $q$ and acyclic chromatic number $c$ has track number at most $c(2q)^{c-1}$
    \end{theorem}
    \begin{proof}
	\begin{itemize}
		\pause
		    \only<-8> {%
		    \item Use the $q$ queues for a $q$-edge-coloring.%
		    }%
		    \only<9-> {%
			\alert<9>{%
		    \item Partition the edges of each queue into upwards and downwards edges, then use the resulting $2q$ sets as a $2q$-edge-coloring.%
		    }%
		    }%
		\pause
	    \item Apply the refinement lemma to an acyclic $c$-coloring and the above edge \alt<9->{\alert<9>{$2q$}}{$q$}-coloring.
		\pause
	    \item The tracks are the refined colors, each ordered like in the queue layout
	\end{itemize}
	\only<9->{
	    Do the ``tracks'' actually form a track layout?
	    \begin{itemize}
		\item<10-> Yes (we fixed them!)
	\end{itemize}
	    }
    \end{proof}
    \pause
    \only<-8>{
	Do the ``tracks'' actually form a track layout? \pause 
	\begin{itemize}
	    \item NO!
		\pause
	    \item But we can fix them
    \end{itemize}
	    }
\end{frame}

\begin{frame}{Outline of the proof}
    \foreach \i in {6, ..., 7} {
	\only<+>{
	    \begin{center}
		\still[\i]{static/outline}
	    \end{center}
	}
    }
\end{frame}

\begin{frame}
    \begin{theorem}[Grünbaum 1972]
	Outerplanar graphs have acyclic chromatic number at most three
    \end{theorem}
    \pause
    \begin{proof}
	Let $G$ be an outerplanar graph. \pause WLoG assume $G$ is a triangulated polygon. Recursively color the graph as follows:\pause
    \begin{tikzpicture}[overlay]
	\fill[black] (-30, -30) rectangle (30, -0.5);
    \end{tikzpicture}
    \end{proof}
\end{frame}

\begin{frame}
    \begin{theorem}[Grünbaum 1972]
	Planar graphs have acyclic chromatic number at most nine.
    \end{theorem}
    \pause
    \begin{proof}
	Let $G$ be a planar graph. \pause WLOG assume $G$ is connected. \pause We will construct a $9$-coloring and show that it does not contain a bichromatic cycle \pause (and that it is proper).
	\begin{itemize}
		\pause
	    \item Choose some \textit{center} $v_0$.
		\pause
	    \item We define the \emph{shell} $S_i = \left\{v \in G | d(v, v_0)=i\right\}$. 
		\pause
	    \item Use a different set of three colors for the shells $S_{3k}$, $S_{3k+1}$, $S_{3k+2}$
		\pause
	    \item Observe: Edges of $G$ are within a single shell or between two adjacent shells.
		\pause
	    \item We will now color each shell $S_i$ such that it cannot be the innermost shell of a bichromatic cycle.
	\end{itemize}
	\phantom\qedhere
    \end{proof}
\end{frame}

\begin{frame}[plain]{Coloring the shells}
    \begin{task}
	Color the shell $S_i$ such that it cannot be the innermost shell of a bichromatic cycle.
    \end{task}
    \begin{tikzpicture}[overlay]
	\fill[black] (-30, -30) rectangle (30, 0);
    \end{tikzpicture}
\end{frame}

\begin{frame}{Concluding the proof}
    \begin{theorem}[Grünbaum 1972]
	Planar graphs have acyclic chromatic number at most nine.
    \end{theorem}
    \pause
    \begin{proof}[Proof (cont.)]
	We have shown that, using our coloring, no bichromatic cycle exists. \pause It is easy to see that the coloring is proper. \pause Therefore, we have found an acyclic $9$-coloring.
    \end{proof}
\end{frame}

\begin{frame}{Better bounds}
    \begin{theorem}[Borodin 1979]
	Planar graphs have acyclic chromatic number at most five.
    \end{theorem}
    \pause
    \begin{proof}
	Same as our proof for outerplanar graphs (i.e. reducing any graph to another graph with less vertices)\pause, except with 450 cases.\only<+(+1)>{\tiny Details left as exercise.}
    \end{proof}
    \pause
    \begin{theorem}[Grünbaum 1972]
	Some planar graphs have acyclic chromatic number five.
    \end{theorem}
\end{frame}

\begin{frame}{Summary}
    \foreach \i in {8, ..., 9} {
	\only<+>{
	    \begin{center}
		\still[\i]{static/outline}
	    \end{center}
	}
    }
\end{frame}

\begin{frame}{Best known upper bounds}
    \foreach \i in {10} {
	\only<+>{
	    \begin{center}
		\still[\i]{static/outline}
	    \end{center}
	}
    }
\end{frame}

\begin{frame}{Best known lower bound of queue number}
    \foreach \i in {11} {
	\only<+>{
	    \begin{center}
		\still[\i]{static/outline}
	    \end{center}
	}
    }
\end{frame}

% \begin{frame}{Summary}
    % \begin{table}
	% \centering
% 
% \begin{tabular}{lllll}
                                  % &   & Acyclic   &      & 3D      \\
                                  % &  Queue &   Chromatic &  Track    &   Drawing    \\
                                  % & Number  &   Number &   Number   &   Volume    \\
% Proven now                        & 42           & \pause9                        &\pause $9\cdot10^{13}$ &\pause $8\cdot 10^{27}\cdot n$ \pause \\
% Proven anytime                    & 42           &\pause 5                        &\pause $3\cdot10^7$    &\pause $4\cdot10^{11}\cdot n$ \pause  \\
% Best lower bound of qn &\pause 4            &\pause 5                        &\pause $1280$          &\pause $1.5\cdot10^6\cdot n$   
% \end{tabular}
    % \end{table}
% \end{frame}

\begin{frame}{Further results}
    \foreach \i in {11,12,13,14} {
	\only<+>{
	    \begin{center}
		\still[\i]{static/outline}
	    \end{center}
	}
    }
\end{frame}

% \begin{frame}{Further results}
    % \begin{itemize}
	% \item Some of the steps we've taken also work in reverse:
	    % \begin{itemize}
		% \item Queue number is bounded by track number.
		% \item If a graph has a $\mathcal O(1) \times \mathcal O(1) \times \mathcal O(n)$ drawing, then it has track number $\mathcal O(1)$
	    % \end{itemize}
	% \item Acyclic chromatic number is also bounded by queue number
	    % \begin{itemize}
		% \item In particular, any proper minor-closed family has 3D drawings with linear volume
	    % \end{itemize}
    % \end{itemize}
    % \end{frame}

    \begin{frame}{Open problems}
    \begin{itemize}
	    \pause
	\item Can we improve the upper bound on
	    \begin{itemize}
		    \pause
		\item the volume of $n$-vertex 3D drawings?
		    \pause
		\item the track number?
		    \pause
		\item the queue number?
	    \end{itemize}
	\item Can we construct linear-volume drawings with 
	    \pause
	    \begin{itemize}
		\item low aspect ratio?
		    \pause
		\item adjacent vertices being close?
		    \pause
		\item other additional properties?
	    \end{itemize}
    \end{itemize}
\end{frame}

\bgroup
\setbeamercolor{background canvas}{bg=black}
\begin{frame}[plain]{}
\end{frame}
\egroup
\end{document}
